    \documentclass[11pt,a4paper]{article}
\usepackage[utf8]{inputenc}
\usepackage[czech]{babel} 
\usepackage{lmodern}
\usepackage[margin=2.5cm, headheight=14pt]{geometry}
\usepackage{fancyhdr}         
\pagestyle{fancy}
\fancyhf{}
\rfoot{\thepage}
\renewcommand{\headrulewidth}{0pt}
\usepackage{amsmath, amsfonts}
\usepackage[most]{tcolorbox}
\usepackage{xcolor}
\usepackage{amssymb}
\usepackage{hyperref}
\usepackage{graphicx}
\usepackage{float}
\usepackage{array}
\usepackage{booktabs}
\usepackage{algorithm}
\usepackage{algpseudocode}
\usepackage{tocloft}
\usepackage{minted}
\usemintedstyle{one-dark}
\usepackage[most]{tcolorbox}
\usepackage{tikz}
\usetikzlibrary{calc}

\setlength{\cftsecnumwidth}{2em}

\DeclareMathOperator{\inDFS}{in}
\DeclareMathOperator{\outDFS}{out}
\DeclareMathOperator{\short}{d}
\DeclareMathOperator{\route}{D}
\DeclareMathOperator{\wei}{w}
\DeclareMathOperator{\score}{f}


\definecolor{modra}{RGB}{66, 153, 255}
\definecolor{cervena}{RGB}{232, 25, 11}
\definecolor{zelena}{RGB}{16, 150, 22}
\definecolor{oranzova}{RGB}{232, 160, 26}

\newtcolorbox{modrybox}[1]{
  colback=gray!5!white,       % Pozadí těla (velmi světle šedá)
  colframe=modra,            % Barva rámečku a hlavičky
  coltitle=white,             % Barva textu v hlavičce
  fonttitle=\bfseries,        % Tučný text v hlavičce
  title=#1,                   % Text v hlavičce předaný jako parametr
  arc=3mm,                    % Zaoblení rohů
  boxrule=1pt,                % Tloušťka vnějšího rámečku
  enhanced,
}

\newtcolorbox{cervenybox}[1]{
  colback=gray!5!white,       % Pozadí těla (velmi světle šedá)
  colframe=cervena,            % Barva rámečku a hlavičky
  coltitle=white,             % Barva textu v hlavičce
  fonttitle=\bfseries,        % Tučný text v hlavičce
  title=#1,                   % Text v hlavičce předaný jako parametr
  arc=3mm,                    % Zaoblení rohů
  boxrule=1pt,                % Tloušťka vnějšího rámečku
  enhanced,
}

\newtcolorbox{zelenybox}[1]{
  colback=gray!5!white,       % Pozadí těla (velmi světle šedá)
  colframe=zelena,            % Barva rámečku a hlavičky
  coltitle=white,             % Barva textu v hlavičce
  fonttitle=\bfseries,        % Tučný text v hlavičce
  title=#1,                   % Text v hlavičce předaný jako parametr
  arc=3mm,                    % Zaoblení rohů
  boxrule=1pt,                % Tloušťka vnějšího rámečku
  enhanced,
}

\newtcolorbox{oranzovybox}[1]{
  colback=gray!5!white,       % Pozadí těla (velmi světle šedá)
  colframe=oranzova,            % Barva rámečku a hlavičky
  coltitle=white,             % Barva textu v hlavičce
  fonttitle=\bfseries,        % Tučný text v hlavičce
  title=#1,                   % Text v hlavičce předaný jako parametr
  arc=3mm,                    % Zaoblení rohů
  boxrule=1pt,                % Tloušťka vnějšího rámečku
  enhanced,
}


\begin{document}

\noindent {\small Počítačové sítě a systémy}

\vspace{0.4em}

\noindent {\Huge\bfseries Simulátor radiace} 

\vspace{0.3em}

\noindent {\Large Dokumentace a rozbor projektu}

\vspace{0.7em}

\noindent \textbf{Autor:} Matěj Pospíšil

\noindent \textbf{Datum:} \today

\vspace{1.2em}
\hrule
\vspace{1.7em}

\tableofcontents

\newpage

\section{Úvod}
\subsection{Představení a ambice}
Mým vybraným projektem je přístroj simulující radiační záření v podobě systému vysílač-přijímač určený pro airsoft či jiné larp akce. Rád bych tímto projektem došel k takové verzi, která bude již alespoň základně využitelná v praxi. Mým cílem je mimo jiné snížit výrobní náklady na minimum, protože přijímačů pro reálnou praxi potřebuji nejméně 10.
\subsection{Funkčnost a plán}
Hráči po lese rozsejí vysílače a každý z nich obdrží přijímač jako PDA. PDA bude neustále skenovat pro vybrané vysílače a vždy si vybere ten nejsilnější signál. Podle jeho síly bude hráči ubírat herní životy a podle MAC adresy vysílače bude určovat momentální lokaci, v níž se hráč nachází. Každé PDA bude mít svůj displej ukazující počet životů a intenzitu radiace a také NFC čtečku, která umožní hráči doplnit životy, vyléčit radiaci či se proti radiaci \uv{obrnit}. Celý tento systém bude zabalený v malé 3D tištěné voděodolné krabičce.

\section{Potřebné komponenty}
\subsection{Mikrokontrolér}
Jako mikrokontrolér jsem zvolil \textit{ESP32-WROOM-32U}, protože má v sobě již zabudovaný WiFi vysílač/přijímač a také tento model umožňuje možnou a snadnou instalaci pomocné antény. Pomocná anténa bude důležitým prvkem, protože stromy, keře a listí velmi silně pohlcují rádiové vlny. Bez případné antény by byl rozsah vysílání maličký. Jako pomocnou anténu do PDA jsem zvolil samolepicí WiFi anténu \textit{GSM 2.4GHz, 4dBi}, kterou pak jen přilepím zevnitř na stranu krabičky. Pro samotný vysílač jsem zvolil anténu \textit{2.4GHz 3dBi lomená s kabelem IPEX U.FL SMA všesměrová}, která zvýší pole působení mnohonásobně. 

\noindent Vzhledem k tomu, že tento modul umí jak vysílat tak přijímat, bude celý proces výroby jednotný a budu moct využívat naprosto stejného vývojového prostředí. Modul je také kompatibilní s Arduino, což mou práci jen usnadní.

\subsection{NFC čtečka}
Jako NFC čtečku jsem zvolil \textit{RFID RC522}, protože je kompatibilní se zvoleným modulem (komunikuje přes SPI sběrnici). Běžně se tato čtečka používá pro různé aplikace v systémech kontroly, ale pro můj projekt bude stačit, když bude načítat NFC tagy.  

\subsection{Displej}
Jako displej jsem zvolil \textit{0.96" OLED displej (I2C)}. Je malý a kompaktní, má malou spotřebu (protože je OLED) a má vysoké rozlišení (128x64 bodů). Dokonce lze zakoupit variantu podporující žlutou barvu, což se pro radiační tematiku hodí. 

\subsection{Napájení}
Vybrat vhodné napájení je ten největší problém. Je zapotřebí obstarat takovou baterii, aby byla schopna celé zapojení utáhnout a hned se nevybít. Proto jsem si vybral klasickou baterii \textit{18650 2500mAh}. Tato baterie by dle očekávání měla vydržet napájet celý přístroj po dobu cca 15 hodin. Je také malá, lehká, kompaktní a jednoduše se zastrčí do držáku.

\noindent Aby bylo možné tuto baterii využít, musím ji nějakým způsobem zavézt do modulu a ohlídat podvybití. K tomuto mi poslouží nabíjecí modul \textit{TP4056}. Baterka se k němu jen připájí a od té doby bude její stav pod kontrolou. Je ovšem potřeba pořídit ten správný modul, protože na trhu jsou dva. Ten správný musí obsahovat \textbf{4} pájecí plošky (B+,B-,OUT+,OUT-).

\noindent Baterie se kvůli časové úspoře nebudou nabíjet, ale vyměňovat. Zařízení bude hráče upozorňovat o stavu baterie a na hráče bude muset být apelováno, aby stav baterie hlídali a vyměnili ji včas. Modul \textit{TP4056} sice chrání před podvybitím, ale je mnohem bezpečnější baterii vyměnit ještě před jeho zásahem.

\section{Komponenty v prototypu}
Můj táta už o takovýhle projekt v minulosti usiloval. Avšak to celé skončilo u nákupu hromady různých dílků. Požádal jsem ho o to, aby mi je svěřil. Prošel jsem si je všechny jednotlivě a rozhodl se, zda lze nějaký z mých naplánovaných dílků vyměnit. 

\noindent Zde popíši, jaké díly vyměním za jaké: 

\subsection{Mikrokontroléry}
Původním plánem bylo využít \textit{ESP32-WROOM-32U}. Místo něj se využiji: 
\begin{itemize}
    \item \textbf{1x ESP32 S3 s 2.4G WiFi modulem} jakožto vysílač (je totiž možné nainstalovat anténu)
    \item \textbf{2x ESP-32S s převodníkem CH340C} jakožto jeden jako vysílač a druhý jako přijímač
\end{itemize}
\noindent Bohužel jeden vysílač bude v prototypu slabší co se týče síly vysílaného signálu. Potřebuji ale mít dva vysílače, abych mohl kontrovat automatickou funkčnost přepnutí přijímaného signálu. 

\noindent Co se týče antén ty se také nachází v balíčku překvapení. Bohužel mi není známý jejich původ. Jsou kompatibilní pouze s \textit{ESP32 S3}.

\subsection{NFC čtečky}
Výhoda pro NFC čtečky je, že v tátovým balení se nachází čtyři úplně stejné čtečky jaké jsem si naplánoval. Tedy \textit{RFID RC522}.

\subsection{Displeje}
Další výhodou je, že v balení od táty se nachází i můj požadovaný displej. Mám ale bohužel jen jeden, takže bude potřeba zbývající displeje dokoupit. Tento displej nemá žluté podsvícení.

\noindent V balení se také nachází další displej, který ale umí zobrazovat pouze znaky na mřížce 16x2. Je také velký, proto není vhodný pro můj projekt.

\subsection{Co je potřeba dokoupit}
Zde je seznam všech dílů, které musím pro první prototypy externě dokoupit:
\begin{itemize}
    \item \textbf{2x} \textit{0.96" OLED displej (I2C)} žluté barvy
    \item \textbf{3x} baterie \textit{18650 2500mAh}
    \item \textbf{3x} modul \textit{TP4056}
    \item \textbf{1x} nepájivé pole (pro prototyp)
\end{itemize}

\subsection{Kontstrukce prototypu}
Prototyp bude konstruován ze součástek zmíněných výše v této kapitole. 
Úplně první \textbf{testovací} prototyp bude realizován na nepájivým poli. 
To je, protože je nebezpečné hned vše k sobě pájet, když ještě není jasné, zda bude systém fungovat. 
\end{document}